\documentclass[10pt,a4paper]{amsart}
\usepackage[margin=19mm]{geometry}
\usepackage{amsmath,amssymb,mathtools}
\usepackage[scr=rsfs,cal=euler]{mathalfa}
\usepackage{xfrac}
\usepackage[unicode,hidelinks,bookmarks=false]{hyperref}
\newcommand{\ie}{i.e.\ }
\newcommand{\cf}{cf.\ }
\newcommand{\resp}{respectively\ }
\newcommand{\Subsection}{\subsection}
\providecommand{\nobreakdash}{\nobreak\hbox{-}}
\setlength{\emergencystretch}{2em}
\multlinegap0pt
\usepackage{bm}
\usepackage{bbm}
\usepackage{dsfont}
\usepackage{xspace}
\usepackage{cancel}
\usepackage{mathtools}
\usepackage{amsthm}
\makeatletter
\@ifundefined{thm}{\newtheorem{thm}{Thm}}{}
\@ifundefined{lem}{\newtheorem{lem}[thm]{Lemma}}{}
\makeatother
\title[Non-existence of Lyapunov exponents in the Newhouse domain]{Non-existence of Lyapunov exponents in the Newhouse domain}
\author{Shin Kiriki, Xiaolong Li, Yushi Nakano, Teruhiko Soma}
\date{Source: arXiv:2604.10913v1 (CC BY 4.0)}
\begin{document}
\maketitle

\noindent\emph{Source and licence.} Non-existence of Lyapunov exponents in the Newhouse domain. Version arXiv:2604.10913v1 (CC BY 4.0), distributed under the Creative Commons Attribution 4.0 International licence, as recorded in the arXiv licence metadata for this exact version and shown on the abstract page https://arxiv.org/abs/2604.10913v1. Full rights notice: licenses/arxiv-2604.10913-CC-BY-4.0.txt.\par\smallskip

\noindent\textit{Editorial scope.} Counted unit: the Lem whose source label is \texttt{lem:250404a} in the article (source0.tex lines 1929--1968, source label \texttt{lem:250404a}), delivered together with the article's own complete original proof of it (source0.tex lines 1970--2208, one \texttt{proof} environment of 239 lines). No mathematics is written, repaired or re-derived: statement and proof are the article's own text line for line. Required context retained uncounted: eq:20251018b (lines 1151--1155); eq:k0decay (lines 1181--1184); eq:0412 (lines 1885--1890); eq:1004a (lines 1892--1896). Every definition, display and label of those lines is retained unchanged. Omitted from this package and not redistributed: 1--1150, 1156--1180, 1185--1884, 1891--1891, 1897--1928, 1969--1969, 2209--2448. Nothing inside a retained excerpt is omitted or altered: every retained body is delivered exactly as the article's lines are written, so no omission receipt of any kind accompanies this package. Citations: the retained excerpts carry no citation command at all, so no bibliography entry of the article is transcribed and no reference is redistributed. Reference adaptation: none was needed and none was made; every reference inside the delivered unit resolves inside it, so no unresolved reference or citation is produced. Printed slips kept literal: no printed word, symbol, hypothesis or number of the counted statement or of its proof was altered. Layout and class: the article's own class is \texttt{amsart} with packages including amsmath, amsthm, amscd, amssymb, tikz, mathrsfs, comment, graphicx, enumerate, color, bm, float; the wrapper is the lane's pinned \texttt{amsart} editorial wrapper at 10pt on A4, which restates the theorem-like environments the retained text needs and adds the article's own macro definitions that the retained text needs (none), with extra wrapper packages (bm, bbm, dsfont, xspace, cancel, mathtools). The delivered numbering is the unit's own. Assets: no retained span contains a figure environment, a graphics-inclusion command, a PDF-page inclusion, a table or a TikZ picture; no image file is copied into this package, the package declares no asset dependency and no image file is redistributed with it. Licence: version arXiv:2604.10913v1 of this article is distributed under the Creative Commons Attribution 4.0 International licence, as recorded in the arXiv licence metadata for this exact version and shown on the abstract page https://arxiv.org/abs/2604.10913v1. A full rights notice is delivered at licenses/arxiv-2604.10913-CC-BY-4.0.txt. Characters: every retained excerpt is pure ASCII, so the delivered engine, XeLaTeX with its default Latin Modern text font, reports no missing character.
    \begin{equation}\label{eq:20251018b}
       \max\{\xi\Lambda_1^{2n_{k_0}^T},
      (\lambda_s\lambda_u^{-1})^{n_{k_0}^H}\Lambda_1^{2n_{k_0}^T}
      \}\le 10^{-4}.
    \end{equation}

    \begin{equation}\label{eq:k0decay}
        \frac{\Lambda_{s}^{(m)}}{\Lambda_{u}^{(m)}}
        \le \frac{1}{12\cdot 3^{m} C_{f,m}^{2}\Lambda_1^{2n_{k_0+m}^T}}.
    \end{equation}

\begin{equation}\label{eq:0412}
    A_l=\begin{bmatrix}
        b_l^{11}\lambda_s^{n_{k_0+l}^H} & b_l^{12}\lambda_u^{n_{k_0+l}^H}\\
        b_l^{21}\lambda_s^{n_{k_0+l}^H} & b_l^{22}\lambda_s^{n_{k_0+l}^H}
    \end{bmatrix},
\end{equation}

\begin{equation}\label{eq:1004a}
    \max_{ij}|b_l^{ij}|\le\Lambda_1^{n_{k_0+l}^T},\quad
    |b_l^{22}|\le\xi\Lambda_1^{n_{k_0+l}^T},\quad
    \min\{|b_l^{12}|,|b_l^{21}|\}\ge\Lambda_1^{-n_{k_0+l}^T}.
\end{equation}

\begin{lem}\label{lem:250404a}
For every odd integer $m\ge 1$, 
\begin{equation}\label{eq:1221-odd}
\left(1-\left(\frac{1}{3}+\frac{1}{3^2}+\cdots +\frac{1}{3^{m}}\right)\right)C_{f,m}^{-1}
\le |C_{ij}^m|
\le
\left(1+\frac{1}{3}+\frac{1}{3^2}+\cdots +\frac{1}{3^{m}}\right)C_{f,m}
\end{equation}
for each $ij\in\{12,21\}$, and
\begin{equation}\label{eq:1221-odd-ratio}
\max\left\{
\frac{|C_{11}^m|}{|C_{12}^m|},
\frac{|C_{22}^m|}{|C_{21}^m|}
\right\}
\le
\frac{1}{3}+\frac{1}{3^2}+\cdots +\frac{1}{3^{m}}.
\end{equation}

For every even integer $m\ge 1$, 
\begin{equation}\label{eq:1221-even}
\left(1-\left(\frac{1}{3}+\frac{1}{3^2}+\cdots +\frac{1}{3^{m}}\right)\right)C_{f,m}^{-1}
\le |C_{ij}^m|
\le
\left(1+\frac{1}{3}+\frac{1}{3^2}+\cdots +\frac{1}{3^{m}}\right)C_{f,m}
\end{equation}
for each $ij\in\{11,22\}$, and
\begin{equation}\label{eq:1221-even-ratio}
\max\left\{
\frac{|C_{12}^m|}{|C_{11}^m|},
\frac{|C_{21}^m|}{|C_{22}^m|}
\right\}
\le
\frac{1}{3}+\frac{1}{3^2}+\cdots +\frac{1}{3^{m}}.
\end{equation}

In particular, for every $m\ge 1$ and $i,j\in\{1,2\}$,
\[
|C_{ij}^m|<2C_{f,m}.
\]
\end{lem}

\begin{proof}
For $m=1$, we have
\[
A^{(1)}=A_0=
\begin{bmatrix}
b_0^{11}\lambda_s^{n_{k_0}^H} & b_0^{12}\lambda_u^{n_{k_0}^H}\\
b_0^{21}\lambda_s^{n_{k_0}^H} & b_0^{22}\lambda_s^{n_{k_0}^H}
\end{bmatrix},
\]
so
\[
C_{11}^1=b_0^{11}\Bigl(\frac{\lambda_s}{\lambda_u}\Bigr)^{n_{k_0}^H},
\qquad
C_{12}^1=b_0^{12},
\qquad
C_{21}^1=b_0^{21},
\qquad
C_{22}^1=b_0^{22}.
\]
By \eqref{eq:1004a},
\[
C_{f,1}^{-1}\le |C_{ij}^1|\le C_{f,1}\quad (ij\in\{12,21\}).
\]
Also,
\[
\frac{|C_{11}^1|}{|C_{12}^1|}
=
\Bigl(\frac{\lambda_s}{\lambda_u}\Bigr)^{n_{k_0}^H}
\left|\frac{b_0^{11}}{b_0^{12}}\right|
\le
\Lambda_1^{2n_{k_0}^T}
\Bigl(\frac{\lambda_s}{\lambda_u}\Bigr)^{n_{k_0}^H}
\le 10^{-4}
\]
by \eqref{eq:1004a} and \eqref{eq:20251018b}. Similarly,
\[
\frac{|C_{22}^1|}{|C_{21}^1|}
=
\left|\frac{b_0^{22}}{b_0^{21}}\right|
\le
\xi \Lambda_1^{2n_{k_0}^T}
\le 10^{-4}
\]
by \eqref{eq:1004a} and \eqref{eq:20251018b}. 
Therefore, \eqref{eq:1221-odd} and \eqref{eq:1221-odd-ratio} hold for $m=1$.

Assume now that the desired estimates hold for $m=l-1$ with an integer $l\ge 2$.
Since
\[
A^{(l)}=A_{l-1}A^{(l-1)},
\quad
\Lambda_u^{(l)}=\lambda_u^{n_{k_0+l-1}^H}\Lambda_s^{(l-1)},
\quad
\Lambda_s^{(l)}=\lambda_s^{n_{k_0+l-1}^H}\Lambda_u^{(l-1)},
\]
it follows from \eqref{eq:0412} that
\begin{align*}
&C_{11}^{l}
=b_{l-1}^{11}C_{11}^{l-1}\frac{\Lambda_s^{(l)}}{\Lambda_u^{(l)}}+b_{l-1}^{12}C_{21}^{l-1},\quad
C_{12}^{l}
=b_{l-1}^{11}C_{12}^{l-1}\frac{\Lambda_s^{(l)}}{\Lambda_u^{(l)}}+b_{l-1}^{12}C_{22}^{l-1},\\
&C_{21}^{l}
=b_{l-1}^{21}C_{11}^{l-1}+b_{l-1}^{22}C_{21}^{l-1}\frac{\Lambda_s^{(l-1)}}{\Lambda_u^{(l-1)}},\quad
C_{22}^{l}
=b_{l-1}^{21}C_{12}^{l-1}+b_{l-1}^{22}C_{22}^{l-1}\frac{\Lambda_s^{(l-1)}}{\Lambda_u^{(l-1)}}.
\end{align*}
Since
$
\frac{1}{3}+\frac{1}{3^2}+\cdots +\frac{1}{3^{l-1}}<\frac12,
$
the inductive hypothesis implies
\[
|C_{ij}^{\,l-1}|<\frac32\,C_{f,l-1}
\]
for all $i,j\in\{1,2\}$.
Indeed, if $l-1$ is odd, then
\[
|C_{ij}^{\,l-1}|
\le
\left(1+\frac{1}{3}+\cdots +\frac{1}{3^{l-1}}\right)C_{f,l-1}
<
\frac32\,C_{f,l-1}
\]
for $ij\in\{12,21\}$,
and
\[
|C_{11}^{\,l-1}|
\le
\left(\frac{1}{3}+\cdots +\frac{1}{3^{l-1}}\right)|C_{12}^{\,l-1}|
<
\frac12\cdot\frac32\,C_{f,l-1}
<
\frac32\,C_{f,l-1},
\]
\[
|C_{22}^{\,l-1}|
\le
\left(\frac{1}{3}+\cdots +\frac{1}{3^{l-1}}\right)|C_{21}^{\,l-1}|
<
\frac12\cdot\frac32\,C_{f,l-1}
<
\frac32\,C_{f,l-1}.
\]
If $l-1$ is even, the same argument with the roles of
$(C_{11}^{\,l-1},C_{22}^{\,l-1})$ and $(C_{12}^{\,l-1},C_{21}^{\,l-1})$
interchanged gives again the claimed estimate.

Hence \eqref{eq:1004a} and \eqref{eq:k0decay} imply
\[
\left|b_{l-1}^{\mu\nu}C_{ij}^{\,l-1}\frac{\Lambda_s^{(l)}}{\Lambda_u^{(l)}}\right|
\le
\Lambda_1^{n_{k_0+l-1}^T}\cdot \frac32\,C_{f,l-1}\cdot
\frac{1}{12\cdot 3^l\,C_{f,l}^2}
=
\frac{1}{8\cdot 3^l}\,C_{f,l}^{-1}
\]
for all relevant indices, because
$
C_{f,l}=\Lambda_1^{n_{k_0+l-1}^T}C_{f,l-1}.
$
Moreover, \eqref{eq:20251018b} implies $\xi<1/3$, and therefore
\begin{align*}
\left|b_{l-1}^{22}C_{ij}^{\,l-1}\frac{\Lambda_s^{(l-1)}}{\Lambda_u^{(l-1)}}\right|
&\le
\xi\Lambda_1^{n_{k_0+l-1}^T}\cdot \frac32\,C_{f,l-1}\cdot
\frac{1}{12\cdot 3^{l-1}\,C_{f,l-1}^2\Lambda_1^{2n_{k_0+l-1}^T}}\\
&=
\frac{\xi}{8\cdot 3^{l-1}\Lambda_1^{n_{k_0+l-1}^T}C_{f,l-1}}\\
&=
\frac{\xi}{8\cdot 3^{l-1}}\,C_{f,l}^{-1}
\le
\frac{1}{8\cdot 3^l}\,C_{f,l}^{-1}
\end{align*}
for all relevant indices.

Suppose first that $l$ is odd. Then $l-1$ is even, so it follows from the above observations that
\begin{align*}
|C_{12}^{l}|
&\le
\frac{1}{8\cdot 3^l}C_{f,l}^{-1}
+\Lambda_1^{n_{k_0+l-1}^T}
\left(1+\frac{1}{3}+\cdots +\frac{1}{3^{l-1}}\right)C_{f,l-1}\\
&\le
\left(1+\frac{1}{3}+\cdots +\frac{1}{3^{l}}\right)C_{f,l},
\end{align*}
because $C_{f,l}^{-1}\le C_{f,l}$ and $C_{f,l}=\Lambda_1^{n_{k_0+l-1}^T}C_{f,l-1}$. Also,
\begin{align*}
|C_{12}^{l}|
&\ge
-\frac{1}{8\cdot 3^l}C_{f,l}^{-1}
+\Lambda_1^{-n_{k_0+l-1}^T}
\left(1-\left(\frac{1}{3}+\cdots +\frac{1}{3^{l-1}}\right)\right)C_{f,l-1}^{-1}\\
&\ge
\left(1-\left(\frac{1}{3}+\cdots +\frac{1}{3^{l}}\right)\right)C_{f,l}^{-1}.
\end{align*}
Exactly the same argument gives
\[
\left(1-\left(\frac{1}{3}+\cdots +\frac{1}{3^{l}}\right)\right)C_{f,l}^{-1}
\le |C_{21}^{l}|
\le
\left(1+\frac{1}{3}+\cdots +\frac{1}{3^{l}}\right)C_{f,l}.
\]
Hence \eqref{eq:1221-odd} holds for $m=l$.

We next prove the ratio estimate \eqref{eq:1221-odd-ratio}.
Let
\[
L_a=b_{l-1}^{11}C_{11}^{\,l-1}\frac{\Lambda_s^{(l)}}{\Lambda_u^{(l)}},
\quad
R_a=b_{l-1}^{12}C_{21}^{\,l-1},\quad
L_b=b_{l-1}^{11}C_{12}^{\,l-1}\frac{\Lambda_s^{(l)}}{\Lambda_u^{(l)}},
\quad
R_b=b_{l-1}^{12}C_{22}^{\,l-1}.
\]
Then
\[
\frac{|R_a|}{|R_b|}
=
\frac{|C_{21}^{\,l-1}|}{|C_{22}^{\,l-1}|}
\le
\frac{1}{3}+\frac{1}{3^2}+\cdots +\frac{1}{3^{l-1}}
\]
by \eqref{eq:1221-even-ratio} in the inductive hypothesis. On the other hand,
\[
|L_a|,\ |L_b|
\le
\frac{1}{8\cdot 3^l}C_{f,l}^{-1},
\qquad
|R_b|
\ge
\left(1-\left(\frac{1}{3}+\cdots +\frac{1}{3^{l-1}}\right)\right)C_{f,l}^{-1}
>
\frac12\,C_{f,l}^{-1},
\]
so
\[
\frac{|L_a|}{|R_b|},\ \frac{|L_b|}{|R_b|}
\le
\frac{1}{4\cdot 3^l}.
\]
Therefore
\begin{align*}
\frac{|C_{11}^{l}|}{|C_{12}^{l}|}
=
\frac{|L_a+R_a|}{|L_b+R_b|}
\le
\frac{|L_a|+|R_a|}{|R_b|-|L_b|}
\le
\frac{\frac{1}{4\cdot 3^l}
+\left(\frac{1}{3}+\cdots +\frac{1}{3^{l-1}}\right)}
{1-\frac{1}{4\cdot 3^l}}
\le
\frac{1}{3}+\frac{1}{3^2}+\cdots +\frac{1}{3^{l}}
\end{align*}
since
$
\frac{1}{4\cdot 3^l}
(
1+\frac{1}{3}+\cdots +\frac{1}{3^{l-1}}+\frac{1}{3^l}
)
\le \frac{1}{3^l}.
$
The proof of
$
\frac{|C_{22}^{l}|}{|C_{21}^{l}|}
\le
\frac{1}{3}+\frac{1}{3^2}+\cdots +\frac{1}{3^{l}}
$
is entirely similar. Thus \eqref{eq:1221-odd-ratio}
holds for $m=l$.

Suppose next that $l$ is even. Then noticing that $l-1$ is odd, 
and repeating the same argument, with the roles of
$(C_{11}^{\,l-1},C_{22}^{\,l-1})$ and $(C_{12}^{\,l-1},C_{21}^{\,l-1})$
interchanged, we obtain
\eqref{eq:1221-even} and \eqref{eq:1221-even-ratio}
for $m=l$.
This completes the induction and the proof.
\end{proof}

\end{document}
